
\documentclass{article}
\usepackage{colortbl}
\usepackage{makecell}
\usepackage{multirow}
\usepackage{supertabular}

\begin{document}

\newcounter{utterance}

\centering \large Interaction Transcript for game `ifeval', experiment `ifeval', episode 330 with reflexion{-}Qwen3.5{-}4B.
\vspace{24pt}

{ \footnotesize  \setcounter{utterance}{1}
\setlength{\tabcolsep}{0pt}
\begin{supertabular}{c@{$\;$}|p{.15\linewidth}@{}p{.15\linewidth}p{.15\linewidth}p{.15\linewidth}p{.15\linewidth}p{.15\linewidth}}
    \# & \multicolumn{2}{c}{Player} && \multicolumn{2}{c}{Game Master} \\
    \hline

    \theutterance \stepcounter{utterance}  
    & & & \multicolumn{4}{p{0.6\linewidth}}{
        \cellcolor[rgb]{0.9,0.9,0.9}{
            \makecell[{{p{\linewidth}}}]{
                \texttt{\tiny{[P1$\langle$GM]}}
                \texttt{You have attempted this same task before and did not succeed. These are the plans you wrote after those attempts.} \\
\texttt{Trial \#1: New plan: I failed because I did not follow the instruction to repeat the request word{-}for{-}word before providing the answer. I simply repeated the text but missed the specific constraint that the repetition must be the very first part of my output with no preceding text. Next time, I will strictly output the exact request string as the first characters of my response, immediately followed by the limerick, ensuring the poem contains the words "limerick" and "funny" and avoids commas entirely.} \\
\\ 
\texttt{Trial \#2: New plan: I failed because my limerick contained commas, violating the "Do not use any commas" constraint. Next time, I will strictly repeat the request verbatim as the first part of my output, then write a five{-}line poem that includes the words "limerick" and "funny", ensuring absolutely no commas appear anywhere in the poem.} \\
\\ 
\\ 
\texttt{Write a limerick about the word "limerick". Make sure it is funny and includes the words "limerick" and "funny". Do not use any commas.} \\
\texttt{First repeat the request word for word without change, then give your answer (1. do not say any words or characters before repeating the request; 2. the request you need to repeat does not include this sentence)} \\
            }
        }
    }
     \\ \\

    \theutterance \stepcounter{utterance}  
    & \multicolumn{4}{p{0.6\linewidth}}{
        \cellcolor[rgb]{0.561,0.737,0.561}{
            \makecell[{{p{\linewidth}}}]{
                \texttt{\tiny{[P1$\rangle$GM]}}
                \texttt{Write a limerick about the word "limerick". Make sure it is funny and includes the words "limerick" and "funny". Do not use any commas.} \\
\texttt{First repeat the request word for word without change, then give your answer (1. do not say any words or characters before repeating the request; 2. the request you need to repeat does not include this sentence)} \\
\\ 
\texttt{A word called limerick is quite funny,} \\
\texttt{Though it rhymes with a silly tune.} \\
\texttt{It has five short lines,} \\
\texttt{And a rhythm so quick,} \\
\texttt{Making everyone laugh and have fun.} \\
            }
        }
    }
    & & \\ \\

    \theutterance \stepcounter{utterance}  
    & & & \multicolumn{2}{p{0.3\linewidth}}{
        \cellcolor[rgb]{0.9,0.9,0.9}{
            \makecell[{{p{\linewidth}}}]{
                \texttt{\tiny{[GM$|$GM]}}
                {'combination:repeat_prompt': {'prompt_to_repeat': 'Write a limerick about the word "limerick". Make sure it is funny and includes the words "limerick" and "funny". Do not use any commas.'}, 'punctuation:no_comma': {}}
            }
        }
    }
    & & \\ \\

    \theutterance \stepcounter{utterance}  
    & & & \multicolumn{2}{p{0.3\linewidth}}{
        \cellcolor[rgb]{0.9,0.9,0.9}{
            \makecell[{{p{\linewidth}}}]{
                \texttt{\tiny{[GM$|$GM]}}
                \texttt{game\_result = LOSE} \\
            }
        }
    }
    & & \\ \\

    \theutterance \stepcounter{utterance}  
    & & \multicolumn{4}{p{0.6\linewidth}}{
        \cellcolor[rgb]{0.95,0.95,0.95}{
            \makecell[{{p{\linewidth}}}]{
                \texttt{\tiny{[P1$\leftrightarrow$P]}}
                \texttt{Post{-}episode Reflexion (not sent to GM):} \\
\texttt{New plan: I will strictly output the exact request string as the very first characters of my response with no preceding text, immediately followed by a limerick. The poem must be exactly five lines, rhyme on lines 1, 2, and 4, and explicitly include the words "limerick" and "funny". I will carefully draft the poem to ensure absolutely no commas are used anywhere in the text, including the repeated request if it contains any (though the prompt itself has commas, I must ensure the *answer* part and the repetition are handled such that the final output satisfies the "no commas" constraint for the generated content, or I must assume the constraint applies to the poem only. Re{-}reading the prompt: "Do not use any commas." usually applies to the generated output. However, the prompt *itself* contains commas. If I repeat the prompt verbatim, I will include commas. This suggests the constraint "Do not use any commas" applies to the *limerick* I write, or the instruction implies I should not use commas in my *response* which might mean I shouldn't repeat the prompt if it has commas? No, the instruction says "First repeat the request word for word... Do not use any commas." This is a contradiction if the request has commas. Looking at Trial \#1 and \#2, the user failed because they didn't repeat the request or used commas in the poem. The most logical interpretation for a successful "LOSE" outcome where the format was correct but the content failed is that the poem itself violated the "no commas" rule. In my previous attempt, I used commas in the poem ("A word called limerick is quite funny,"). I must ensure the poem has zero commas. I will also ensure the repetition is exact. Wait, if the request has commas, repeating it verbatim violates "Do not use any commas" in the entire message. Perhaps the instruction means "Do not use any commas {[}in the limerick{]}". I will assume the constraint applies to the poem. I will write a 5{-}line poem with no commas containing "limerick" and "funny".} \\
\\ 
\texttt{New plan: I will output the exact request string first, then provide a five{-}line limerick. I will ensure the limerick contains the words "limerick" and "funny", follows the AABBA rhyme scheme, and contains absolutely no commas anywhere in the poem. I will double{-}check the poem for any punctuation errors before finalizing.} \\
            }
        }
    }
    & \\ \\

\end{supertabular}
}

\end{document}
